\section*{Bab \ref{ch:b1:rock-forming-minerals} --- Mineral pembentuk batuan}

\begin{solution}{exo:b1:rock-forming-minerals:1}
\ce{SiO4^{4-}}: \omterm{def:b1:rock-forming-minerals:silicate-classes}{nesosilikat}, $s = 0$. \ce{Si2O7^{6-}}: \omterm{def:b1:rock-forming-minerals:silicate-classes}{sorosilikat},
$s = 1$. \ce{Si6O18^{12-}}: \omterm{def:b1:rock-forming-minerals:silicate-classes}{siklosilikat}, $s = 2$. \ce{Si2O5^{2-}}:
\omterm{def:b1:rock-forming-minerals:silicate-classes}{filosilikat}, $s = 3$. \ce{SiO2}: \omterm{def:b1:rock-forming-minerals:silicate-classes}{tektosilikat}, $s = 4$. \ce{SiO3^{2-}}
memiliki $4 - s/2 = 3$, sehingga $s = 2$: sebuah cincin (\omterm{def:b1:rock-forming-minerals:silicate-classes}{siklosilikat}) atau rantai tunggal
(\omterm{def:b1:rock-forming-minerals:silicate-classes}{inosilikat}); rasio saja tidak dapat membedakannya.
\end{solution}

\begin{solution}{exo:b1:rock-forming-minerals:2}
Forsterit: $2\times24.305 + 28.085 + 4\times15.999 =
\qty{140.691}{g.mol^{-1}}$; fayalit: $2\times55.845 + 28.085 +
4\times15.999 = \qty{203.771}{g.mol^{-1}}$. Volume molar:
$140.691/3.275 = \qty{42.96}{cm^3.mol^{-1}}$ dan $203.771/4.392 =
\qty{46.40}{cm^3.mol^{-1}}$.
\end{solution}

\begin{solution}{exo:b1:rock-forming-minerals:3}
$(0.78 - 0.72)/0.72 = +\qty{8.3}{\percent}$ untuk \ce{Fe^{2+}};
$-\qty{4.2}{\percent}$ untuk \ce{Ni^{2+}}; $+\qty{15.3}{\percent}$ untuk
\ce{Mn^{2+}}; $+\qty{38.9}{\percent}$ untuk \ce{Ca^{2+}}. Besi dan nikel, di bawah \qty{15}{\percent}, menggantikan magnesium secara bebas; mangan, sedikit di atasnya, terbatas; kalsium, di atas \qty{30}{\percent}, hampir tidak masuk. Keempatnya memiliki muatan yang sama dengan magnesium.
\end{solution}

\begin{solution}{exo:b1:rock-forming-minerals:4}
$(3.51 - 2.16)/2.16 = \qty{62.5}{\percent}$ dari grafit ke intan;
$(2.93 - 2.71)/2.71 = \qty{8.1}{\percent}$ dari kalsit ke aragonit. \omterm{def:b1:rock-forming-minerals:polymorph}{Polimorf} yang lebih padat, intan dan aragonit, adalah bentuk tekanan tinggi.
\end{solution}

\begin{solution}{exo:b1:rock-forming-minerals:5}
Dengan $x = 0.30$ dalam $\mathrm{Na}_{1-x}\mathrm{Ca}_x\mathrm{Al}_{1+x}\mathrm{Si}_{3-x}\mathrm{O}_8$:
$\mathrm{(Na_{0.7}Ca_{0.3})(Al_{1.3}Si_{2.7})O_8}$. Muatan:
$0.7 + 2\times0.3 + 3\times1.3 + 4\times2.7 = 0.7 + 0.6 + 3.9 + 10.8 = +16$,
dan $8\times(-2) = -16$.
\end{solution}

\begin{solution}{exo:b1:rock-forming-minerals:6}
Satu kelompok berulang memiliki dua tetrahedron yang berbagi tiga sudut (yang puncaknya menghubungkan kedua rantai) dan dua yang berbagi dua. Keempat tetrahedron memiliki $4\times4 = 16$ sudut, dari yang $3 + 3 + 2 + 2 = 10$ dibagi, yaitu
$10/2 = 5$ oksigen yang dibagi; $16 - 10 = 6$ lainnya tidak dibagi. Total
$5 + 6 = 11$ oksigen untuk 4 silikon: \ce{Si4O11}, dengan muatan
$4\times4 - 11\times2 = -6$. Proposisi memberikan hal yang sama dengan
$s = 5/2$: $4 - 5/4 = 11/4$ oksigen per silikon, muatan $5/2 - 4 = -3/2$ per
silikon.
\end{solution}

\begin{solution}{exo:b1:rock-forming-minerals:7}
Mol dalam \qty{100}{g}: \ce{SiO2} $40.20/60.083 = 0.66907$; \ce{MgO}
$47.10/40.304 = 1.16862$; \ce{FeO} $12.20/71.844 = 0.16981$; \ce{NiO}
$0.40/74.692 = 0.00536$. Oksigen: $2\times0.66907 + 1.16862 + 0.16981 +
0.00536 = 2.68194$; $f = 4/2.68194 = 1.49146$. Per empat oksigen: Si
\num{0.998}, Mg \num{1.743}, Fe \num{0.253}, Ni \num{0.008}. Situs oktahedral memuat $1.743 + 0.253 + 0.008 = 2.004$, situs tetrahedral
\num{0.998}: analisis yang baik, \ce{(Mg_{1.743}Fe_{0.253}Ni_{0.008})SiO4}.
$\mathrm{Fo} = 100\times1.743/(1.743 + 0.253) = 87.3$. Faktor $f$
mengalikan Mg dan Fe secara sama, sehingga ia saling menghapus dalam perbandingan:
$\mathrm{Fo} = 100\,(47.10/40.304)/(47.10/40.304 + 12.20/71.844) =
100\times1.16862/1.33843 = 87.3$.
\end{solution}

\begin{solution}{exo:b1:rock-forming-minerals:8}
Positif: $1 + 2\times3 + 3\times4 + 3 = +22$; negatif:
$10\times(-2) + 2\times(-1) = -22$. Lembaran dengan komposisi \ce{Si4O10} akan membawa $-4$ (\cref{prop:b1:rock-forming-minerals:si-o}, $s = 3$); menggantikan satu silikon dari empat dengan aluminium menambah satu muatan negatif, sehingga \ce{(Si3Al)O10} membawa $-5$. Dua aluminium oktahedral dengan dua hidroksil, $6 - 2 = +4$, menyeimbangkan empat dari lima; ion kalium di antara lembaran menyeimbangkan yang kelima.
\end{solution}

\begin{solution}{exo:b1:rock-forming-minerals:9}
(a) $M = (140.691 + 203.771)/2 = \qty{172.231}{g.mol^{-1}}$, $V = (42.96 +
46.40)/2 = \qty{44.68}{cm^3.mol^{-1}}$, $\rho = 172.231/44.68 =
\qty{3.855}{g.cm^{-3}}$; rata-rata massa jenis adalah $(3.275 + 4.392)/2 =
\qty{3.834}{g.cm^{-3}}$, lebih rendah \qty{0.021}{g.cm^{-3}}. (b) Selesaikan
$203.771 + x(140.691 - 203.771) = 3.600\,\bigl(46.40 + x(42.96 - 46.40)\bigr)$:
$203.771 - 63.080x = 167.040 - 12.384x$, sehingga $x = 36.731/50.696 = 0.7245$: \omterm{def:b1:rock-forming-minerals:groups}{olivin} itu adalah $\text{Fo}_{72}$ (hingga persen terdekat).
\end{solution}

\begin{solution}{exo:b1:rock-forming-minerals:10}
Dengan $M_i = \rho_i V_i$,
\[
  \rho(x) - \bigl(x\rho_1 + (1 - x)\rho_2\bigr)
  = \frac{x\rho_1V_1 + (1 - x)\rho_2V_2 - \bigl(x\rho_1 + (1 - x)\rho_2\bigr)\bigl(xV_1 + (1 - x)V_2\bigr)}{V(x)} .
\]
Dengan menjabarkan hasil kali, pembilangnya adalah $x\rho_1V_1 - x^2\rho_1V_1 +
(1 - x)\rho_2V_2 - (1 - x)^2\rho_2V_2 - x(1 - x)(\rho_1V_2 + \rho_2V_1)$, dan
$x - x^2 = x(1 - x)$, $(1 - x) - (1 - x)^2 = x(1 - x)$, sehingga ia adalah
$x(1 - x)(\rho_1V_1 + \rho_2V_2 - \rho_1V_2 - \rho_2V_1) =
x(1 - x)(\rho_1 - \rho_2)(V_1 - V_2)$. Untuk \omterm{def:b1:rock-forming-minerals:groups}{olivin} (1 = forsterit) kedua
selisih itu negatif, $3.275 - 4.392 = -1.117$ dan $42.96 - 46.40 =
-3.44$, sehingga hasil kalinya positif dan massa jenisnya berada di atas garis lurus. Pada $\text{Fo}_{50}$: $0.25\times1.117\times3.44/44.68 =
\qty{0.0215}{g.cm^{-3}}$, yang sesuai dengan \cref{exo:b1:rock-forming-minerals:9} hingga pembulatan volume molarnya.
\end{solution}

\begin{solution}{exo:b1:rock-forming-minerals:11}
Total \qty{99.60}{\percent}. Mol oksida: \ce{SiO2} $0.90042$,
\ce{Al2O3} $0.28344$, \ce{CaO} $0.19973$, \ce{Na2O} $0.08229$,
\ce{K2O} $0.00318$. Oksigen: $1.80084 + 0.85033 + 0.19973 + 0.08229 +
0.00318 = 2.93636$; $f = 8/2.93636 = 2.72446$. Kation (dua kali mol untuk
\ce{Al2O3}, \ce{Na2O}, \ce{K2O}): Si \num{2.453}, Al \num{1.544}, Ca
\num{0.544}, Na \num{0.448}, K \num{0.017}. Situs tetrahedral:
$2.453 + 1.544 = 3.997$; situs besar: $0.544 + 0.448 + 0.017 = 1.009$. Maka
$\mathrm{Al} - 1 = 0.544$, sama dengan Ca dalam pembulatan: setiap kalsium masuk bersama satu aluminium. $\mathrm{An} = 100\times0.544/1.009 = 53.9$: sebuah
labradorit, $\text{An}_{54}$.
\end{solution}

\begin{solution}{exo:b1:rock-forming-minerals:12}
(a) $M(\ce{Fe3O4}) = 3\times55.845 + 4\times15.999 = \qty{231.531}{g.mol^{-1}}$.
Tiga besinya, dilaporkan sebagai \ce{FeO}, berbobot $3\times71.844 = \qty{215.532}{g}$,
yaitu $215.532/231.531 = \qty{93.09}{\percent}$; oksigen yang terikat pada
\ce{Fe^{3+}}, di luar yang diperhitungkan \ce{FeO}, hilang dari total.
(b) $93.09/71.844 = 1.29572$ mol \ce{FeO} dan oksigen; $f =
4/1.29572 = 3.08708$; empat besi per empat oksigen, \ce{Fe4O4}. Magnetit memiliki
tiga kation per empat oksigen. Oksigennya kurang dihitung: setiap
\ce{Fe^{3+}} memegang satu setengah oksigen, bukan satu seperti yang diasumsikan \ce{FeO}. (c) Pada tiga kation, Fe $= 3$
dengan $a$ \ce{Fe^{2+}} dan $b$ \ce{Fe^{3+}}: $a + b = 3$ dan $2a + 3b = 8$,
sehingga $b = 2$, $a = 1$: $\mathrm{Fe^{2+}Fe^{3+}_2O_4}$. (d) Satu \ce{FeO} dan
satu \ce{Fe2O3} per unit rumus: $71.844/231.531 = \qty{31.03}{\percent}$
\ce{FeO} dan $159.687/231.531 = \qty{68.97}{\percent}$ \ce{Fe2O3}, total
\qty{100.00}{\percent}.
\end{solution}

\begin{solution}{pb:b1:rock-forming-minerals:1}
\textbf{1.} \ce{SiO2} $28.085 + 2\times15.999 = 60.083$; \ce{FeO} 71.844;
\ce{MnO} 70.937; \ce{MgO} 40.304; \ce{NiO} 74.692; \ce{CaO} 56.077
(\si{g.mol^{-1}}).

\textbf{2.} $39.50 + 16.53 + 0.23 + 43.32 + 0.20 + 0.22 =
\qty{100.00}{\percent}$. Total \qty{97}{\percent} akan menyarankan sebuah komponen yang tidak dianalisis (air, atau oksigen yang terikat pada \ce{Fe^{3+}} yang dilaporkan sebagai \ce{FeO}), atau analisis yang buruk.

\textbf{3.} \ce{SiO2} $39.50/60.083 = 0.65742$; \ce{FeO} 0.23008; \ce{MnO}
0.00324; \ce{MgO} 1.07483; \ce{NiO} 0.00268; \ce{CaO} 0.00392.

\textbf{4.} $2\times0.65742 = 1.31485$ untuk \ce{SiO2}, angka yang sama dengan
soal 3 untuk yang lain; jumlah $1.31485 + 0.23008 + 0.00324 + 1.07483 +
0.00268 + 0.00392 = 2.62960$.

\textbf{5.} Setiap oksida memiliki satu kation: angka-angka soal 3.

\textbf{6.} $f = 4/2.62960 = 1.52114$.

\textbf{7.} Si $0.65742\times1.52114 = 1.000$; Fe 0.350; Mn 0.005; Mg 1.635;
Ni 0.004; Ca 0.006.

\textbf{8.} $1.635 + 0.350 + 0.005 + 0.004 + 0.006 = 2.000$ untuk situs
oktahedral dan 1.000 untuk situs tetrahedral: tepat 2 dan
1 yang ideal.

\textbf{9.} \ce{(Mg_{1.635}Fe_{0.350}Mn_{0.005}Ni_{0.004}Ca_{0.006})SiO4}.

\textbf{10.} $2\times2.000 + 4\times1.000 = +8.000$ terhadap $4\times(-2) =
-8$. Setiap kation dihitung dengan muatan oksidanya, sehingga muatan positifnya dua kali oksigen secara konstruksi: pemeriksaan ini tidak mungkin gagal. Jumlah situs soal 8 bisa gagal.

\textbf{11.} Nikel berbeda dari magnesium sebesar
$(0.69 - 0.72)/0.72 = -\qty{4}{\percent}$, jauh di dalam \qty{15}{\percent}:
ia menggantikan secara bebas. Kalsium $(1.00 - 0.72)/0.72 = \qty{39}{\percent}$
lebih besar, di atas \qty{30}{\percent}: ia hampir tidak masuk.

\textbf{12.} Besi yang dihitung sebagai \ce{FeO} ketika sebagiannya adalah \ce{Fe^{3+}} memberikan terlalu sedikit oksigen per kation, sehingga lebih dari tiga kation pada empat oksigen
(\cref{exo:b1:rock-forming-minerals:12}). Jumlahnya $1.000 + 2.000 = 3.000$: tidak ada
\ce{Fe^{3+}} yang berarti.

\textbf{13.} $\mathrm{Fo} = 100\times1.635/(1.635 + 0.350) = 163.5/1.985 =
82.4$.

\textbf{14.} $\mathrm{Mg\#} = 82.4$: semua besi adalah \ce{Fe^{2+}}
(soal 12), sehingga $\mathrm{Fe^{2+}} = \mathrm{Fe}$ dan kedua definisi berimpit.

\textbf{15.} Dengan $\text{Fo}_{82.4}\text{Fa}_{17.6}$, $x = 0.824$; besi mengisi
\qty{17.6}{\percent} situs oktahedral dalam model biner
($0.350/2.000 = \qty{17.5}{\percent}$ dengan unsur-unsur kecil).

\textbf{16.} Faktor itu mengalikan Mg dan Fe secara sama dan saling menghapus:
$\mathrm{Fo} = 100\times1.07483/(1.07483 + 0.23008) = 82.4$.

\textbf{17.} $40.5/40.304 = 1.00486$ dan $20.0/71.844 = 0.27838$;
$\mathrm{Fo} = 100\times1.00486/1.28324 = 78.3$. \omterm{def:b1:rock-forming-minerals:groups}{Olivin} yang kemudian lebih kaya
besi.

\textbf{18.} $M = 140.691$ dan \qty{203.771}{g.mol^{-1}};
$V = 140.691/3.275 = 42.96$ dan $203.771/4.392 =
\qty{46.40}{cm^3.mol^{-1}}$.

\textbf{19.} $M = 0.824\times140.691 + 0.176\times203.771 = 115.929 +
35.864 = \qty{151.793}{g.mol^{-1}}$; $V = 0.824\times42.96 +
0.176\times46.40 = 35.399 + 8.166 = \qty{43.565}{cm^3.mol^{-1}}$.

\textbf{20.} $\rho = 151.793/43.565 = \qty{3.484}{g.cm^{-3}}$.

\textbf{21.} $0.824\times3.275 + 0.176\times4.392 = 2.699 + 0.773 =
\qty{3.472}{g.cm^{-3}}$, \qty{0.012}{g.cm^{-3}} (\qty{0.3}{\percent}) di bawah
nilai campuran ideal: dekat, seperti yang diprediksi
\cref{exo:b1:rock-forming-minerals:10}.

\textbf{22.} $0.12\times\qty{3484}{kg.m^{-3}}\times\qty{1}{m^3} =
\qty{418}{kg}$.

\textbf{23.} $418/2950 = \qty{14.2}{\percent}$: \omterm{def:b1:rock-forming-minerals:groups}{olivin} yang padat membentuk porsi massa yang lebih besar daripada porsi volumenya.

\textbf{24.} \omterm{def:b1:rock-forming-minerals:groups}{Olivin} itu adalah $\text{Fo}_{82}$: kandungan forsteritnya adalah
\textbf{82 persen}.
\end{solution}
